\documentclass[11pt]{article}
\usepackage{climbingguide}

\begin{document}

% -------------------
% INDEX PAGE REPLICA
% -------------------
\CGIndexCorner
\CGIndexBackground

\vspace*{2mm}
{\setlength{\leftskip}{8mm}
\CGBaseFont\normalsize
\textbf{ATENÇÃO!} \dotfill 10 \\
Ética local \dotfill 11 \\
Como usar o guia \dotfill 12 \\

\textbf{Zona Boqueirão} \dotfill 30 \\
\hspace*{1em}Mocozado \dotfill 33 \\
\hspace*{1em}Mocozinho \dotfill 34 \\
\hspace*{1em}Mocó \dotfill 36 \par
}

\clearpage

% -------------------
% ZONE COVER REPLICA
% -------------------
\CGZoneHeader{example-image-a}{ZONA BOQUEIRÃO}

\begin{multicols}{2}
    {\CGTitleFont\LARGE\color{brand_color} Setores}\par\vspace{1pt}
    {\CGBaseFont\normalsize Mocozado, Mocozinho e Mocó.}\par\vspace{6pt}
    A Zona Boqueirão é composta por três setores distintos. A caminhada leva 20 minutos.
    \columnbreak

    {\CGTitleFont\LARGE\color{brand_color} Betas da zona}\par\vspace{1pt}
    O horário mais recomendado é à tarde. O uso de repelente é indispensável.\par\vspace{6pt}
    
    \CGObsBox{\textbf{\color{brand_color}OBS.} É comum encontrar vespas nas fendas das vias. Elas não apresentam comportamento agressivo.}
\end{multicols}

\clearpage

% -------------------
% SECTOR HEADER REPLICA
% -------------------
\CGSectorHeader{sector_yellow}{MOCOZINHO}

\begin{multicols}{2}
    Diferente do Mocó, o setor Mocozinho é mais acessível. A base é confortável e possui vias muito convidativas.
    \columnbreak

    % Auto-indexed and auto-colored routes
    \CGRoute{CAFÉ SEM TABACO}{6sup}{15m}{5+2}{Allan Carlos}
    \CGRoute{EQUILÍBRIO MALUCO}{7c}{15m}{4+2}{Vinicin}
    \CGRoute{NÃO SUFOQUE O ARTISTA}{5sup}{15m}{4+2}{Allan Carlos}
\end{multicols}

\end{document}
